\section{Билет 16. Вынужденные колебания. Уравнение вынужденных колебаний. Описание поведение системы при вынужденных колебаниях. Резонанс. Вывод формулы резонансной частоты. Резонансные кривые. Применение резонанса в технике, отрицательные и полезные свойства.}

\begin{definition}
\textbf{Вынужденными колебаниями} называются колебания, возникающие в системе под действием внешней периодически изменяющейся силы $F(t) = F_0 \cos(\omega t)$. Уравнение движения имеет вид:
\begin{equation}
    m\ddot{x} + r\dot{x} + kx = F_0 \cos(\omega t).
\end{equation}
Разделив на $m$ и введя обозначения $2\beta = r/m$, $\omega_0^2 = k/m$, $f_0 = F_0/m$, получаем неоднородное дифференциальное уравнение:
\begin{equation}
    \ddot{x} + 2\beta \dot{x} + \omega_0^2 x = f_0 \cos(\omega t). \label{eq:forced_16}
\end{equation}
\end{definition}

\begin{theorem}[Поведение системы и установившиеся колебания]
Общее решение уравнения \eqref{eq:forced_16} складывается из затухающего собственного колебания и частного решения. После завершения переходного процесса ($t \gg 1/\beta$) собственные колебания затухают, и система совершает \textbf{установившиеся вынужденные колебания} на частоте внешней силы $\omega$:
\begin{equation}
    x(t) = A \cos(\omega t - \varphi).
\end{equation}
Амплитуда $A$ и сдвиг фазы $\varphi$ определяются параметрами системы и частотой вынуждающей силы:
\begin{equation}
    A = \frac{f_0}{\sqrt{(\omega_0^2 - \omega^2)^2 + 4\beta^2 \omega^2}}, \quad \tan\varphi = \frac{2\beta\omega}{\omega_0^2 - \omega^2}. \label{eq:amplitude_phase}
\end{equation}
\end{theorem}

\begin{property}[Резонансная частота]
\textbf{Резонансом} называется явление резкого возрастания амплитуды вынужденных колебаний при приближении частоты внешней силы к некоторому критическому значению. Для нахождения резонансной частоты $\omega_{\text{рез}}$ найдём минимум знаменателя в формуле для $A$. Возьмём производную по $\omega$ и приравняем к нулю:
\begin{equation}
    \frac{d}{d\omega}\left[ (\omega_0^2 - \omega^2)^2 + 4\beta^2 \omega^2 \right] = 0.
\end{equation}
\begin{equation}
    2(\omega_0^2 - \omega^2)(-2\omega) + 8\beta^2 \omega = 0 \quad \Rightarrow \quad -4\omega(\omega_0^2 - \omega^2 - 2\beta^2) = 0.
\end{equation}
Отбрасывая тривиальный корень $\omega = 0$, получаем:
\begin{equation}
    \omega_{\text{рез}} = \sqrt{\omega_0^2 - 2\beta^2}. \label{eq:res_freq}
\end{equation}
При малом затухании ($\beta \ll \omega_0$) резонансная частота практически совпадает с собственной: $\omega_{\text{рез}} \approx \omega_0$.
\end{property}

\begin{remark}
\textbf{Резонансные кривые.} Зависимость амплитуды $A$ от частоты $\omega$ называется амплитудно-частотной характеристикой. Чем меньше коэффициент затухания $\beta$, тем:
\begin{enumerate}
    \item Выше и уже резонансный пик;
    \item Ближе $\omega_{\text{рез}}$ к $\omega_0$;
    \item Больше добротность системы $Q$.
\end{enumerate}
\textbf{Применение и свойства резонанса:}
\begin{itemize}
    \item \textbf{Полезные:} настройка музыкальных инструментов, радиоприёмники (избирательность контура), МРТ-диагностика, акустические резонаторы, раскачивание качелей.
    \item \textbf{Вредные:} разрушение мостов и зданий при совпадении частоты внешних нагрузок с собственной частотой конструкции, обрыв проводов, вибрации трубопроводов, раскачивание грузов на кранах.
    \item \textbf{Защита:} изменение собственной частоты конструкции, установка виброгасителей, разделение частот, сбивание шага при прохождении мостов.
\end{itemize}
\end{remark}
